\begin{frame}{A silly upper bound}
    \begin{columns}
        \begin{column}{0.5\textwidth}
            
            \begin{theorem}{Theorem}
                \(D(n,k) \leq \binom{n}{k-1}\)
            \end{theorem}
            
            \begin{proof}{}
                In the reconfiguration path, each edge has sets with \(k-1\) vertices in common on both sides. \\
                \pause
                The path is minimal, so they don't repeat. Thus the number of edges is at most \(\binom{n}{k-1}\).
            \end{proof}
            
        \end{column}
    
        \begin{column}{0.5\textwidth}
        
            \begin{figure}
                \centering
                \begin{tikzpicture}[scale=1.5]
        
    \node[circle, draw=black, fill=white, minimum size=0.6cm, line width=1.5pt] (v1) at (1, 2) {$1$};
    \node[circle, draw=black, fill=white, minimum size=0.6cm, line width=1.5pt] (v2) at (3, 2) {$2$};
    \node[circle, draw=black, fill=white, minimum size=0.6cm, line width=1.5pt] (v3) at (0, 1) {$3$};
    \node[circle, draw=black, fill=white, minimum size=0.6cm, line width=1.5pt] (v4) at (2, 1) {$4$};
    \node[circle, draw=black, fill=white, minimum size=0.6cm, line width=1.5pt] (v5) at (4, 1) {$5$};
    \node[circle, draw=black, fill=white, minimum size=0.6cm, line width=1.5pt] (v6) at (1, 0) {$6$};
    \node[circle, draw=black, fill=white, minimum size=0.6cm, line width=1.5pt] (v7) at (3, 0) {$7$};
            
            \draw[line width=1.5pt, dashed] (v1) -- (v2);
            \draw[line width=1.5pt, dashed] (v1) -- (v3);
            \draw[line width=1.5pt, dashed] (v1) -- (v4);
            \draw[line width=1.5pt, dashed] (v3) -- (v6);
            \draw[line width=1.5pt, dashed] (v1) -- (v6);
            \draw[line width=1.5pt, dashed] (v4) -- (v6);
            \draw[line width=1.5pt, dashed] (v2) -- (v4);
            \draw[line width=1.5pt, dashed] (v2) -- (v5);
            \draw[line width=1.5pt, dashed] (v2) -- (v7);
            \draw[line width=1.5pt, dashed] (v4) -- (v7);
            \draw[line width=1.5pt, dashed] (v5) -- (v7);
            \draw[line width=1.5pt, dashed] (v6) -- (v7);

        
        \node[circle, draw=cfblue, fill=cfred, minimum size=0.3cm, line width=1.5pt] (r1) at (1.3, 1) {};
        % Bottom left triangle face center: (2/3, 1/3)
        \node[circle, draw=cfblue, fill=cfred, minimum size=0.3cm, line width=1.5pt] (r2) at (0.5, 1) {};
        % Center quadrilateral face center: (2, 1)
        \node[circle, draw=cfblue, fill=cfblue, minimum size=0.3cm, line width=1.5pt] (r3) at (2, 1.6) {};
        % Right quadrilateral face center: (3, 1)
        \node[circle, draw=cfblue, fill=cfred, minimum size=0.3cm, line width=1.5pt] (r4) at (2.6, 1) {};
        % Bottom right triangle face center: (10/3, 1/3)
        \node[circle, draw=cfblue, fill=cfred, minimum size=0.3cm, line width=1.5pt] (r5) at (3.5, 1) {};
        % Bottom center quadrilateral face center: (2, 0.5)
        \node[circle, draw=cfblue, fill=cfred, minimum size=0.3cm, line width=1.5pt] (r6) at (2, 0.3) {};


    \draw[line width=1.5pt, cfblue] (r1) -- (r2);
    \draw[line width=1.5pt, cfblue] (r1) -- (r3);
    \draw[line width=1.5pt, cfblue] (r1) -- (r6);
    \draw[line width=1.5pt, cfblue] (r3) -- (r4);
    \draw[line width=1.5pt, cfblue] (r4) -- (r5);
    \draw[line width=1.5pt, cfblue] (r4) -- (r6);


    \node[circle, draw=black, fill=cfred, minimum size=0.6cm, line width=1.5pt] (p0) at (0, -1) {\footnotesize$1,3,6$};
    \node[circle, draw=black, fill=cfred, minimum size=0.6cm, line width=1.5pt] (p1) at (1, -1) {\footnotesize$1,4,6$};
    \node[circle, draw=black, fill=cfred, minimum size=0.6cm, line width=1.5pt] (p2) at (2, -1) {\footnotesize$1,2,4$};
    \node[circle, draw=black, fill=cfred, minimum size=0.6cm, line width=1.5pt] (p3) at (3, -1) {\footnotesize$2,4,7$};
    \node[circle, draw=black, fill=cfred, minimum size=0.6cm, line width=1.5pt] (p4) at (4, -1) {\footnotesize$2,5,7$};

\draw[line width=1.5pt] 
  (p0) -- node[midway, above, font=\footnotesize] {$1,6$} (p1);
\draw[line width=1.5pt] 
  (p1) -- node[midway, above, font=\footnotesize] {$1,4$} (p2);
\draw[line width=1.5pt] 
  (p2) -- node[midway, above, font=\footnotesize] {$2,4$} (p3);
\draw[line width=1.5pt] 
  (p3) -- node[midway, above, font=\footnotesize] {$2,7$} (p4);


\end{tikzpicture}

            \end{figure}
            
        \end{column}
    \end{columns}
\end{frame}

%--------------------------------------------

\begin{frame}
    \begin{theorem}{Theorem}
        \(D(n,2)=n-2\)
    \end{theorem}
    \begin{proof}{Proof}
        Independent sets are exactly non edges of \(G\). Thus \(D(n-2)\leq \mathrm{diam}(\overline{G})-1\leq n-2\). \pause Consider complement of a path:
    \begin{figure}
        \centering
                \begin{tikzpicture}[scale=1]
        
            \node[circle, draw=black, fill=white, minimum size=0.4cm, line width=1.5pt] (v1) at (0, 2) {};
            \node[circle, draw=black, fill=white, minimum size=0.4cm, line width=1.5pt] (v2) at (2, 2) {};
            \node[circle, draw=black, fill=white, minimum size=0.4cm, line width=1.5pt] (v3) at (4, 2) {};
            \node[circle, draw=black, fill=white, minimum size=0.4cm, line width=1.5pt] (v4) at (6, 2) {};
            \node[circle, draw=black, fill=white, minimum size=0.4cm, line width=1.5pt] (v5) at (8, 2) {};
            \node[circle, draw=black, fill=white, minimum size=0.4cm, line width=1.5pt] (v6) at (10, 2) {};
            \node[circle, draw=black, fill=white, minimum size=0.4cm, line width=1.5pt] (v7) at (12, 2) {};
            \draw[line width=1.5pt, dashed] (v1) -- (v2);
            \draw[line width=1.5pt, dashed] (v2) -- (v3);
            \draw[line width=1.5pt, dashed] (v3) -- (v4);
            \draw[line width=1.5pt, dashed] (v4) -- (v5);
            \draw[line width=1.5pt, dashed] (v5) -- (v6);
            \draw[line width=1.5pt, dashed] (v6) -- (v7);
            \onslide<3->{
                \draw[cfblue, line width=1.5pt] ($(v1)!0.5!(v2)$) ellipse (1.3cm and 0.6cm);
    \draw[cfblue, line width=1.5pt] ($(v2)!0.5!(v3)$) ellipse (1.3cm and 0.6cm);
    \draw[cfblue, line width=1.5pt] ($(v3)!0.5!(v4)$) ellipse (1.3cm and 0.6cm);
    \draw[cfblue, line width=1.5pt] ($(v4)!0.5!(v5)$) ellipse (1.3cm and 0.6cm);
    \draw[cfblue, line width=1.5pt] ($(v5)!0.5!(v6)$) ellipse (1.3cm and 0.6cm);
    \draw[cfblue, line width=1.5pt] ($(v6)!0.5!(v7)$) ellipse (1.3cm and 0.6cm);
    
            \node[circle, draw=black, fill=cfblue, minimum size=0.4cm, line width=1.5pt] (v8) at (1, 1) {};
            \node[circle, draw=black, fill=cfblue, minimum size=0.4cm, line width=1.5pt] (v9) at (3, 1) {};
            \node[circle, draw=black, fill=cfblue, minimum size=0.4cm, line width=1.5pt] (v10) at (5, 1) {};
            \node[circle, draw=black, fill=cfblue, minimum size=0.4cm, line width=1.5pt] (v11) at (7, 1) {};
            \node[circle, draw=black, fill=cfblue, minimum size=0.4cm, line width=1.5pt] (v12) at (9, 1) {};
            \node[circle, draw=black, fill=cfblue, minimum size=0.4cm, line width=1.5pt] (v13) at (11, 1) {};
            \draw[line width=1.5pt] (v8) -- (v9);
            \draw[line width=1.5pt] (v9) -- (v10);
            \draw[line width=1.5pt] (v10) -- (v11);
            \draw[line width=1.5pt] (v11) -- (v12);
            \draw[line width=1.5pt] (v12) -- (v13);
            }
  
        \end{tikzpicture}    
    \end{figure}
    
    \pause
    
    Thus \(D(n,2)=n-2\)
    
    \end{proof}
\end{frame}

%--------------------------------------------

\begin{frame}
    \begin{columns}

    \begin{column}{0.6\textwidth}
    
        \begin{theorem}{Theorem}
        For \( k \geq 3 \), we have
        \[
        D(n,k) = o\bigl(n^{k-1}\bigr).
        \]
        \end{theorem}
        
        \begin{heading}{Proof idea}
        
            \begin{itemize}
                \item each independent set on diameter path creates a \(k\) vertex, \(k-1\) uniform hyperclique
                \item too many hypercliques \(\to\) shortcut exists
                \item less \(\to\) apply the hypergraph removal lemma
            \end{itemize}
            
        \end{heading}
    
    \end{column}
    \begin{column}{0.4\textwidth}
    
        \begin{figure}
            \centering
            
%                \only<1>{\begin{tikzpicture}[scale=1.3]
%             % graph
%    \node[circle, draw=black, fill=white, minimum size=0.6cm, line width=1.5pt] (v0) at (2, 3) {$0$};
%     \node[circle, draw=black, fill=white, minimum size=0.6cm, line width=1.5pt] (v1) at (1, 2) {$1$};
%     \node[circle, draw=black, fill=white, minimum size=0.6cm, line width=1.5pt] (v2) at (3, 2) {$2$};
%     \node[circle, draw=black, fill=white, minimum size=0.6cm, line width=1.5pt] (v3) at (0, 1) {$3$};
%     \node[circle, draw=black, fill=white, minimum size=0.6cm, line width=1.5pt] (v4) at (2, 1) {$4$};
%     \node[circle, draw=black, fill=white, minimum size=0.6cm, line width=1.5pt] (v5) at (4, 1) {$5$};
%     \node[circle, draw=black, fill=white, minimum size=0.6cm, line width=1.5pt] (v6) at (1, 0) {$6$};
%     \node[circle, draw=black, fill=white, minimum size=0.6cm, line width=1.5pt] (v7) at (3, 0) {$7$};
            
%             \draw[line width=1.5pt] (v0) -- (v1);
%             \draw[line width=1.5pt] (v1) -- (v2);
%             \draw[line width=1.5pt] (v0) -- (v2);
%             \draw[line width=1.5pt] (v1) -- (v3);
%             \draw[line width=1.5pt] (v1) -- (v4);
%             \draw[line width=1.5pt] (v3) -- (v6);
%             \draw[line width=1.5pt] (v1) -- (v6);
%             \draw[line width=1.5pt] (v4) -- (v6);
%             \draw[line width=1.5pt] (v2) -- (v4);
%             \draw[line width=1.5pt] (v2) -- (v5);
%             \draw[line width=1.5pt] (v2) -- (v7);
%             \draw[line width=1.5pt] (v4) -- (v7);
%             \draw[line width=1.5pt] (v5) -- (v7);
%             \draw[line width=1.5pt] (v6) -- (v7);


%         \node[circle, draw=cfblue, fill=cfblue, minimum size=0.3cm, line width=1.5pt] (r0) at (2, 2.3) {};
        
%         \node[circle, draw=cfblue, fill=cfblue, minimum size=0.3cm, line width=1.5pt] (r1) at (1.3, 1) {};
%         % Bottom left triangle face center: (2/3, 1/3)
%         \node[circle, draw=cfblue, fill=cfblue, minimum size=0.3cm, line width=1.5pt] (r2) at (0.5, 1) {};
%         % Center quadrilateral face center: (2, 1)
%         \node[circle, draw=cfblue, fill=cfblue, minimum size=0.3cm, line width=1.5pt] (r3) at (2, 1.6) {};
%         % Right quadrilateral face center: (3, 1)
%         \node[circle, draw=cfblue, fill=cfblue, minimum size=0.3cm, line width=1.5pt] (r4) at (2.6, 1) {};
%         % Bottom right triangle face center: (10/3, 1/3)
%         \node[circle, draw=cfblue, fill=cfblue, minimum size=0.3cm, line width=1.5pt] (r5) at (3.5, 1) {};
%         % Bottom center quadrilateral face center: (2, 0.5)
%         \node[circle, draw=cfblue, fill=cfblue, minimum size=0.3cm, line width=1.5pt] (r6) at (2, 0.3) {};

%     % Reconfiguration edges
%     \draw[line width=1.5pt, cfblue, dashed] (r0) -- (r3);
%     \draw[line width=1.5pt, cfblue, dashed] (r1) -- (r2);
%     \draw[line width=1.5pt, cfblue, dashed] (r1) -- (r3);
%     \draw[line width=1.5pt, cfblue, dashed] (r1) -- (r6);
%     \draw[line width=1.5pt, cfblue, dashed] (r3) -- (r4);
%     \draw[line width=1.5pt, cfblue, dashed] (r4) -- (r5);
%     \draw[line width=1.5pt, cfblue, dashed] (r4) -- (r6);
    
%             \end{tikzpicture}}
%     \skipindraft[Higlights diameter]{
%     \onslide<2->{
%     \begin{tikzpicture}[scale=1.3]
        
%                 % graph
%    \node[circle, draw=black, fill=white, minimum size=0.6cm, line width=1.5pt] (v0) at (2, 3) {$0$};
%     \node[circle, draw=black, fill=white, minimum size=0.6cm, line width=1.5pt] (v1) at (1, 2) {$1$};
%     \node[circle, draw=black, fill=white, minimum size=0.6cm, line width=1.5pt] (v2) at (3, 2) {$2$};
%     \node[circle, draw=black, fill=white, minimum size=0.6cm, line width=1.5pt] (v3) at (0, 1) {$3$};
%     \node[circle, draw=black, fill=white, minimum size=0.6cm, line width=1.5pt] (v4) at (2, 1) {$4$};
%     \node[circle, draw=black, fill=white, minimum size=0.6cm, line width=1.5pt] (v5) at (4, 1) {$5$};
%     \node[circle, draw=black, fill=white, minimum size=0.6cm, line width=1.5pt] (v6) at (1, 0) {$6$};
%     \node[circle, draw=black, fill=white, minimum size=0.6cm, line width=1.5pt] (v7) at (3, 0) {$7$};
            
%             \draw[line width=1.5pt] (v0) -- (v1);
%             \draw[line width=1.5pt] (v1) -- (v2);
%             \draw[line width=1.5pt] (v0) -- (v2);
%             \draw[line width=1.5pt] (v1) -- (v3);
%             \draw[line width=1.5pt] (v1) -- (v4);
%             \draw[line width=1.5pt] (v3) -- (v6);
%             \draw[line width=1.5pt] (v1) -- (v6);
%             \draw[line width=1.5pt] (v4) -- (v6);
%             \draw[line width=1.5pt] (v2) -- (v4);
%             \draw[line width=1.5pt] (v2) -- (v5);
%             \draw[line width=1.5pt] (v2) -- (v7);
%             \draw[line width=1.5pt] (v4) -- (v7);
%             \draw[line width=1.5pt] (v5) -- (v7);
%             \draw[line width=1.5pt] (v6) -- (v7);


%         \node[circle, draw=cfblue, fill=cfblue, minimum size=0.3cm, line width=1.5pt] (r0) at (2, 2.3) {};
        
%         \node[circle, draw=cfblue, fill=cfred, minimum size=0.3cm, line width=1.5pt] (r1) at (1.3, 1) {};
%         % Bottom left triangle face center: (2/3, 1/3)
%         \node[circle, draw=cfblue, fill=cfred, minimum size=0.3cm, line width=1.5pt] (r2) at (0.5, 1) {};
%         % Center quadrilateral face center: (2, 1)
%         \node[circle, draw=cfblue, fill=cfred, minimum size=0.3cm, line width=1.5pt] (r3) at (2, 1.6) {};
%         % Right quadrilateral face center: (3, 1)
%         \node[circle, draw=cfblue, fill=cfred, minimum size=0.3cm, line width=1.5pt] (r4) at (2.6, 1) {};
%         % Bottom right triangle face center: (10/3, 1/3)
%         \node[circle, draw=cfblue, fill=cfred, minimum size=0.3cm, line width=1.5pt] (r5) at (3.5, 1) {};
%         % Bottom center quadrilateral face center: (2, 0.5)
%         \node[circle, draw=cfblue, fill=cfblue, minimum size=0.3cm, line width=1.5pt] (r6) at (2, 0.3) {};

%     % Reconfiguration edges
%     \draw[line width=1.5pt, cfblue, dashed] (r0) -- (r3);
%     \draw[line width=1.5pt, cfblue, dashed] (r1) -- (r2);
%     \draw[line width=1.5pt, cfblue, dashed] (r1) -- (r3);
%     \draw[line width=1.5pt, cfblue, dashed] (r1) -- (r6);
%     \draw[line width=1.5pt, cfblue, dashed] (r3) -- (r4);
%     \draw[line width=1.5pt, cfblue, dashed] (r4) -- (r5);
%     \draw[line width=1.5pt, cfblue, dashed] (r4) -- (r6);
%         \end{tikzpicture}
% }
% }

\only<2>{
\begin{tikzpicture}[scale=1.3]
        
    \node[circle, draw=black, fill=white, minimum size=0.6cm, line width=1.5pt] (v1) at (1, 2) {$1$};
    \node[circle, draw=black, fill=white, minimum size=0.6cm, line width=1.5pt] (v2) at (3, 2) {$2$};
    \node[circle, draw=black, fill=white, minimum size=0.6cm, line width=1.5pt] (v3) at (0, 1) {$3$};
    \node[circle, draw=black, fill=white, minimum size=0.6cm, line width=1.5pt] (v4) at (2, 1) {$4$};
    \node[circle, draw=black, fill=white, minimum size=0.6cm, line width=1.5pt] (v5) at (4, 1) {$5$};
    \node[circle, draw=black, fill=white, minimum size=0.6cm, line width=1.5pt] (v6) at (1, 0) {$6$};
    \node[circle, draw=black, fill=white, minimum size=0.6cm, line width=1.5pt] (v7) at (3, 0) {$7$};
            
            \draw[line width=1.5pt, dashed] (v1) -- (v2);
            \draw[line width=1.5pt, dashed] (v1) -- (v3);
            \draw[line width=1.5pt, dashed] (v1) -- (v4);
            \draw[line width=1.5pt, dashed] (v3) -- (v6);
            \draw[line width=1.5pt, dashed] (v1) -- (v6);
            \draw[line width=1.5pt, dashed] (v4) -- (v6);
            \draw[line width=1.5pt, dashed] (v2) -- (v4);
            \draw[line width=1.5pt, dashed] (v2) -- (v5);
            \draw[line width=1.5pt, dashed] (v2) -- (v7);
            \draw[line width=1.5pt, dashed] (v4) -- (v7);
            \draw[line width=1.5pt, dashed] (v5) -- (v7);
            \draw[line width=1.5pt, dashed] (v6) -- (v7);

        
        \node[circle, draw=cfblue, fill=cfred, minimum size=0.3cm, line width=1.5pt] (r1) at (1.3, 1) {};
        % Bottom left triangle face center: (2/3, 1/3)
        \node[circle, draw=cfblue, fill=cfred, minimum size=0.3cm, line width=1.5pt] (r2) at (0.5, 1) {};
        % Center quadrilateral face center: (2, 1)
        \node[circle, draw=cfblue, fill=cfblue, minimum size=0.3cm, line width=1.5pt] (r3) at (2, 1.6) {};
        % Right quadrilateral face center: (3, 1)
        \node[circle, draw=cfblue, fill=cfred, minimum size=0.3cm, line width=1.5pt] (r4) at (2.6, 1) {};
        % Bottom right triangle face center: (10/3, 1/3)
        \node[circle, draw=cfblue, fill=cfred, minimum size=0.3cm, line width=1.5pt] (r5) at (3.5, 1) {};
        % Bottom center quadrilateral face center: (2, 0.5)
        \node[circle, draw=cfblue, fill=cfred, minimum size=0.3cm, line width=1.5pt] (r6) at (2, 0.3) {};


    \draw[line width=1.5pt, cfblue] (r1) -- (r2);
    \draw[line width=1.5pt, cfblue] (r1) -- (r3);
    \draw[line width=1.5pt, cfblue] (r1) -- (r6);
    \draw[line width=1.5pt, cfblue] (r3) -- (r4);
    \draw[line width=1.5pt, cfblue] (r4) -- (r5);
    \draw[line width=1.5pt, cfblue] (r4) -- (r6);
\end{tikzpicture}
}
        \end{figure}
        
    \end{column}
    \end{columns}
    
    \begin{figure}
        \centering
        \onslide<2>{
\begin{tikzpicture}[scale=1.5]
    \node[circle, draw=black, fill=cfred, minimum size=0.6cm, line width=1.5pt] (p0) at (0, 0) {\footnotesize$1,3,6$};
    \node[circle, draw=black, fill=cfred, minimum size=0.6cm, line width=1.5pt] (p1) at (1, 0) {\footnotesize$1,4,6$};
    \node[circle, draw=black, fill=cfred, minimum size=0.6cm, line width=1.5pt] (p2) at (2, 0) {\footnotesize$1,2,4$};
    \node[circle, draw=black, fill=cfred, minimum size=0.6cm, line width=1.5pt] (p3) at (3, 0) {\footnotesize$2,4,7$};
    \node[circle, draw=black, fill=cfred, minimum size=0.6cm, line width=1.5pt] (p4) at (4, 0) {\footnotesize$2,5,7$};

    \draw[line width=1.5pt] (p0) -- (p1);
    \draw[line width=1.5pt] (p1) -- (p2);
    \draw[line width=1.5pt] (p2) -- (p3);
    \draw[line width=1.5pt] (p3) -- (p4);
\end{tikzpicture}
}
    \end{figure}
\end{frame}

%--------------------------------------------

\begin{frame}

    \begin{columns}
        \begin{column}{0.8\textwidth}
            \onslide<2->{
                We partition the vertices into \textcolor{cfyel}{even \(P_e\)} and \textcolor{cfred}{odd \(P_o\)}.
            }
            
            \\\vspace{1em}
            
            \onslide<3->{
                Intersection of sets in \(P_i\) is at most \(k-2\) thus each edge comes from one independent set.
            }

            \onslide<4->{
                \\\vspace{1em}
                
                \(H_e,H_o\) : \(k-1\) uniform hypergraphs on vertices from \(G\). Edges are independent sets contained in some independent set from \(P_i\).
            }
            \begin{columns}
                \begin{column}{0.55\textwidth}
            \\\vspace{1em}

            \onslide<5->{
                \(K\): \(k-1\) uniform hyperclique of size \(k\).
                
                \\\vspace{1em}
                
                Each \(Z_i\) creates exactly one copy of \(K\)
            }
            \end{column}
            \begin{column}{0.4\textwidth}
                \onslide<4->{
                    \begin{figure}
                        \centering
                        \begin{tikzpicture}[scale=0.8]
    \node[circle, draw=black, fill=white, line width=1.5pt] (p6) at (0, 1) {\footnotesize$6$};
    \node[circle, draw=black, fill=white, line width=1.5pt] (p3) at (0, -1) {\footnotesize$3$}
    \node[circle, draw=black, fill=white, line width=1.5pt] (p1) at (1.5, 0) {\footnotesize$1$}
    \node[circle, draw=black, fill=white, line width=1.5pt] (p2) at (3, 0) {\footnotesize$2$}
    \node[circle, draw=black, fill=white, line width=1.5pt] (p4) at (2.25, 1.5) 
    {\footnotesize$4$}
    \node[circle, draw=black, fill=white, line width=1.5pt] (p5) at (4.5, 1)    {\footnotesize$5$}
    \node[circle, draw=black, fill=white, line width=1.5pt] (p7) at (4.5, -1) 
    {\footnotesize$7$}

    \draw[line width=1.5pt, cfblue] (p6) -- (p1);
    \draw[line width=1.5pt, cfblue] (p6) -- (p3);
    \draw[line width=1.5pt, cfpink] (p5) -- (p7);
    \draw[line width=1.5pt, cfblue] (p3) -- (p1);
    \draw[line width=1.5pt, cfgr] (p1) -- (p4);
    \draw[line width=1.5pt, cfgr] (p2) -- (p4);
    \draw[line width=1.5pt, cfgr] (p1) -- (p2);
    \draw[line width=1.5pt, cfpink] (p2) -- (p5);
    \draw[line width=1.5pt, cfpink] (p2) -- (p7);
\end{tikzpicture}

                    \end{figure}
                }
            \end{column}
            \end{columns}

        \end{column}
        
        \begin{column}{0.2\textwidth}
        
            \begin{figure}
                
\begin{tikzpicture}[scale=1.5]
    \coordinate (p0) at (0, 4);
    \coordinate (p1) at (0, 3);
    \coordinate (p2) at (0, 2);
    \coordinate (p3) at (0, 1);
    \coordinate (p4) at (0, 0);

    \draw[line width=1.5pt] (p0) -- (p1);
    \draw[line width=1.5pt] (p1) -- (p2);
    \draw[line width=1.5pt] (p2) -- (p3);
    \draw[line width=1.5pt] (p3) -- (p4);
    
    \onslide<3->{
        \draw[line width=1.5pt, bend right=80, dashed, cfred] (p0) to node[midway, left] {$1$} (p2);;
        \draw[line width=1.5pt, bend right=80, dashed, cfred] (p2)  to node[midway, left] {$2$} (p4);
    }
    
    \only<1>{
        \node[circle, draw=black, fill=cfred, minimum size=0.6cm, line width=1.5pt, label=right:$Z_1$] at (p0)  {\footnotesize$1,3,6$};
        \node[circle, draw=black, fill=cfred, minimum size=0.6cm, line width=1.5pt, label=right:$Z_2$] at (p1)  {\footnotesize$1,4,6$};
        \node[circle, draw=black, fill=cfred, minimum size=0.6cm, line width=1.5pt, label=right:$Z_3$] at (p2) {\footnotesize$1,2,4$};
        \node[circle, draw=black, fill=cfred, minimum size=0.6cm, line width=1.5pt, label=right:$Z_4$] at (p3) {\footnotesize$2,4,7$};
        \node[circle, draw=black, fill=cfred, minimum size=0.6cm, line width=1.5pt, label=right:$Z_5$] at (p4) {\footnotesize$2,5,7$};
    }

    \only<2-3>{
        \node[circle, draw=black, fill=cfred, minimum size=0.6cm, line width=1.5pt, label=right:$Z_1$] at (p0){$1,3,6$};
        \node[circle, draw=black, fill=cfyel, minimum size=0.6cm, line width=1.5pt, label=right:$Z_2$] at (p1) {$1,4,6$};
        \node[circle, draw=black, fill=cfred, minimum size=0.6cm, line width=1.5pt, label=right:$Z_3$] at (p2) {$1,2,4$};
        \node[circle, draw=black, fill=cfyel, minimum size=0.6cm, line width=1.5pt, label=right:$Z_4$] at (p3) {$2,4,7$};
        \node[circle, draw=black, fill=cfred, minimum size=0.6cm, line width=1.5pt, label=right:$Z_5$] at (p4) {$2,5,7$};
    }

        \only<4->{
        \node[circle, draw=black, fill=cfblue, minimum size=0.6cm, line width=1.5pt, label=right:$Z_1$] at (p0){$1,3,6$};
        \node[circle, draw=black, fill=white, minimum size=0.6cm, line width=1.5pt, label=right:$Z_2$] at (p1) {$1,4,6$};
        \node[circle, draw=black, fill=cfgr, minimum size=0.6cm, line width=1.5pt, label=right:$Z_3$] at (p2) {$1,2,4$};
        \node[circle, draw=black, fill=white, minimum size=0.6cm, line width=1.5pt, label=right:$Z_4$] at (p3) {$2,4,7$};
        \node[circle, draw=black, fill=cfpink, minimum size=0.6cm, line width=1.5pt, label=right:$Z_5$] at (p4) {$2,5,7$};
    }

\end{tikzpicture}

            \end{figure}
            
        \end{column}
    \end{columns}

\end{frame}

%--------------------------------------------

\begin{frame}
    \begin{header}\textcolor{cfred}{
    \textbf{Case 1:}} \(H_i\) contains more than \(n^{k-1}\) copies of \(K\).\\
    Each \(Z_j\in H_i\) adds exactly one copy. There are at most \(n^{k-1}\) sets in \(H_i\), thus there exists copy of \(K\) - \(K'\) such that \(V(K')\notin P_i\).\\
    \vspace{1em}
    \begin{columns}
    \begin{column}{0.53\textwidth}
    In \(K'\) there are 3 hyperedges \(e_1,e_2,e_3\) intersecting on \(k-2\) vertices.\\
    \vspace{1em}
    There exists \(x_1,x_2,x_3 \in V(G)\) such that \(x_i \cup e_i\) are independent sets in \(P_i\).\\
    \vspace{1em}
    \end{column}
    \begin{column}{0.4\textwidth}
        \begin{figure}
            \centering
            \begin{tikzpicture}[scale=0.8]
    \node[circle, draw=black, fill=white, line width=1.5pt] (p6) at (0, 1) {\footnotesize$6$};
    \node[circle, draw=black, fill=white, line width=1.5pt] (p3) at (0, -1) {\footnotesize$3$}
    \node[circle, draw=black, fill=white, line width=1.5pt] (p1) at (1.5, 0) {\footnotesize$1$}
    \node[circle, draw=black, fill=white, line width=1.5pt] (p2) at (3, 0) {\footnotesize$2$}
    \node[circle, draw=black, fill=white, line width=1.5pt] (p4) at (2.25, 1.5) 
    {\footnotesize$4$}
    \node[circle, draw=black, fill=white, line width=1.5pt] (p5) at (4.5, 1)    {\footnotesize$5$}
    \node[circle, draw=black, fill=white, line width=1.5pt] (p7) at (4.5, -1) 
    {\footnotesize$7$}

    \draw[line width=1.5pt, cfblue] (p6) -- (p1);
    \draw[line width=1.5pt, cfblue] (p6) -- (p3);
    \draw[line width=1.5pt, cfpink] (p5) -- (p7);
    \draw[line width=1.5pt, cfblue] (p3) -- (p1);
    \draw[line width=1.5pt, cfgr] (p1) -- (p4);
    \draw[line width=1.5pt, cfgr] (p2) -- (p4);
    \draw[line width=1.5pt, cfgr] (p1) -- (p2);
    \draw[line width=1.5pt, cfpink] (p2) -- (p5);
    \draw[line width=1.5pt, cfpink] (p2) -- (p7);
\end{tikzpicture}

        \end{figure}
    \end{column}
    \end{columns}
    \pause
    But \(e_1 \cup e_2 \cup e_3\) is also independent set in \(G\) of size \(k\) adjacent to \(x_i \cup e_i\). Thus \(x_i \cup e_i\) are at distance 2 from each other, thus contradiction.    
\end{header}
\end{frame}

%--------------------------------------------

\begin{frame}
     \begin{header}\textcolor{cfred}{\textbf{Case 2:}} \(H_i\) contains less than \(n^{k-1}\) copies of \(K\).
     \begin{lemma}{Graph removal lemma}
        \(H\), k-uniform, n-vertex hypergraph\\
        \(K\), r-uniform hypergraph\\
        If \(H\) contains at most \(o(n^k)\) copies of \(K\) then we can eliminate them all by removing at most \(o(n^r)\) hyperedges.
     \end{lemma}
     \pause
     Thus there exists a set of hyperedges \(S\), \(S=|o(n^{k-1}|)\), such that \(H-S\) contains~no~copy~of~\(K\).\\
     \vspace{1em}
     Each hyperedge of \(S\) is contained in exactly one element of \(P_i\), thus \(|P_i|\leq |S| = o(n^{k-1})\). \\
     \[d\leq 2 | P_i|=o(n^{k-1})\].
\end{header}
\end{frame}

%--------------------------------------------


\begin{frame}{Summary: Upper bound}
    \begin{theorem}{Theorem}
        \(D(n,k) =o(n^{k-1})\)
    \end{theorem}
    \textcolor{cfred}{\textbf{Case 1:}} \(H_i\) contains more than \(n^{k-1}\) copies of \(K\).\\
    \(\longrightarrow\) a shortcut exists\\
    \textcolor{cfred}{\textbf{Case 2:}} \(H_i\) contains less than \(n^{k-1}\) copies of \(K\).\\
     \(\longrightarrow\)    Each hyperedge of \(S\) is contained in exactly one element of \(P_i\), thus \(|P_i|\leq |S| = o(n^{k-1})\). Moreover \(d\leq 2 | P_i|=o(n^{k-1})\).
\end{frame}